\begin{proof}

\begin{enumerate}[label={$\langle1\rangle$\arabic*.},leftmargin=*]
\item The universal Picard family constructed in C.1 restricts to a holomorphic family on the proper space $F_c\times P$. Its underlying sheaf is flat over $P$, because it is a line-bundle twist of the fixed locally free restriction $\Omega_Y^1|_{F_c}$. Proper semicontinuity therefore makes $Z_c$ a closed complex analytic subset of $P$.

\item Assume that each $Z_c$ is proper. The complement of their finite union is open dense and has full Lebesgue measure in $P\simeq\C^q$. Statement \ref{cdp:st:D.1.11} supplies a dense set of full measure of line bundles non-torsion on every full fibre. Choose $L$ in the intersection. On every smooth torus fibre, $K_Y|_F$ and the conormal are trivial. The intrinsic cotangent bundle is trivial as well, so the conormal exact sequence expresses the ambient cotangent restriction as an extension of trivial bundles. The line $L|_F$ has zero integral Chern class, since $c_1(L)$ is torsion and $H^2(F,\Z)$ is torsion-free; it is nontrivial by the chosen non-torsion condition. Its cohomology in every degree vanishes, as follows by the unitary Fourier resolution of a topologically trivial line on a torus. The ambient section group is consequently zero at every smooth fibre. At each critical fibre it is zero by the definition of the chosen complement.

\item Serre duality on each Gorenstein Cartier fibre, with $\omega_{F_c}=K_Y|_{F_c}$, and top-degree base change give $R^2f_*(T_Y\otimes L)=0$. Twisting by $f^*\OO(n)$ preserves every fibre restriction. For sufficiently large $n$, it kills $H^1(\PP^1,R^1f_*(T_Y\otimes L)\otimes\OO(n))$. We choose $n$ divisible by the finite order of $c_1(f^*\OO(1))$, so this twist stays in $P$. The two-column Leray sequence gives an actual line bundle in $P$ with zero $H^2$. Proper semicontinuity makes this zero locus dense open. Intersecting it with the dense open endpoint-vanishing loci of C.1 proves \eqref{cdp:eq:C.19}, using $\chi=c_3/2$.

\item Conversely, suppose $H^2$ vanishes on a dense open subset of $P$, and choose $L$ in that subset and in the full-measure set of \ref{cdp:st:D.1.11}. The preceding smooth-fibre argument makes $R^2f_*(T_Y\otimes L)$ a torsion sheaf. Leray gives a surjection
\[
 H^2(Y,T_Y\otimes L)\longrightarrow
 H^0(\PP^1,R^2f_*(T_Y\otimes L)).
\]
Thus that torsion sheaf is zero. Top-degree base change and fibre duality imply $L\notin Z_c$ for every $c$. In particular none of the $Z_c$ equals $P$, proving the equivalence.

\item When $b_3=0$, C.1 gives $H^2(Y,T_Y\otimes L)\ne0$ for every $L$. The equivalence therefore forces some $Z_c$ to be all of $P$. There are only finitely many $c$, and $P$ is an irreducible vector-space torsor; a finite union of proper closed analytic subsets cannot exhaust it. The integral hypothesis makes $P=\Pic(Y)$ by the exponential sequence. This proves the assertion about a fixed fibre. Finally the identifications in A.1--A.3 prove the claimed explicit forms of the fibre tests, without replacing a thickened fibre by its support.

\end{enumerate}
\end{proof}
